% Thesis v3, Chuong 6 Conclusion (danh so, Le 12/11). Viet tu thesis loi V3, so theo phuong an A cua RQ3.
\chapter{Conclusion}\label{ch:conclusion}

\section{Summary of the Findings}\label{sec:c-summary}

\enlargethispage{\baselineskip}\looseness=-1 This thesis examined how far the PASS/FAIL decisions published for C-to-Rust translations are supported by behavioral evidence. It took the released translations of ten systems on the TRACTOR Public Test Corpus together with the decisions recorded for them, reproduced the decisions that rest on running a translation, and compared the behavior of each translation with that of its C program using an independent differential testing harness. The procedure checks the evidence behind each decision directly, and it keeps the decisions without a result in the reported rate.

\looseness=-1 The three research questions received the following answers. For RQ1, \frp{rqFPExeAll} executable and \frp{rqFPLibAll} library translations that had been accepted and were fully re-evaluated have at least one retained divergence from the C program. Every one of these translations diverges on a generated input, and \rqDivTransLayerSExeAllNum{} executable translations also diverge on an original test vector. For RQ2, \frc{rqFFExeAll} executable and \frc{rqFFLibAll} library rejections at the test stage lack supporting evidence among those that reached an outcome, while \rqLrExeAllNum{} executable and \rqLrLibAllNum{} library rejections reached none. For RQ3, the success rate of executables, \rqSRExeAllPct\% under the published decisions, lies between \rqSRLoExeAllPct\% and \srup{Exe}{}\% after reclassification, and that of libraries, \rqSRLibAllPct\% under the published decisions, lies between \rqSRLoLibAllPct\% and \rqSRUpLibAllPct\%. The width of each interval reflects the decisions whose result remains open.

\looseness=-1 On the corpus studied, the evidence that changes the reading of an evaluation therefore lies mainly on the side of acceptances. A translation that passes a fixed test suite can still differ from its C program on inputs that the suite does not contain, whereas the rejections at the test stage that reached an outcome were, with one exception, backed by an observed difference. This supports complementing a reported pass rate with a differential check of the accepted translations on further inputs, and reporting the original population, the decisions actually checked and the causes of attrition together with the rate. These results, together with the re-evaluation procedure that produced them, form the scientific contribution of the thesis. Its technical contribution is a harness for executable programs and library functions that implements the procedure and records every input, observation and adjudication, so that each reported number can be traced; the code and result files are publicly available (Appendix~\ref{app:s-env}).

\enlargethispage{\baselineskip}\looseness=-1 These conclusions hold for the translations, the inputs and the comparison rules that were examined, and the result for rejections reflects a corpus whose expected values are obtained from the C programs. A retained divergence is an observed difference under the oracle of this study: the C program may contain behavior that the sanitizers do not identify, and the absence of a retained divergence does not establish that a translation is equivalent. The divergences that the library controls show to be missed, the non-random attrition and the restriction to one corpus, discussed in Section~\ref{sec:d-threats}, bound how far the results can be extended.

\section{Future Work}\label{sec:c-future}

\looseness=-1 The first priority is to reduce the share of decisions that remain open. Rerunning the heavy cases that were interrupted and completing the missing seeds and records would narrow the interval of RQ3 and show whether the main results hold on a more completely re-evaluated set, and resolving the failure of the sanitizer build would bring the remaining executable rejections into the proportion of RQ2. For libraries, the mutant-based validation showed that some non-equivalent mutants were missed. Further work can develop a guided generation of call states and assess it with independently labeled mutants.

\looseness=-1 A second direction is to cover the kinds of decisions and the conditions of the reference that lie outside the present measurement. Rejections at the build stage need a re-evaluation procedure of their own, which separates failures caused by the translation from failures caused by the environment or the build configuration and defines the corresponding outcomes and denominators. With the C program as the reference, the detection of undefined behavior can be strengthened, and differences that arise from implementation-defined or unspecified behavior can be identified. Finally, applying the same re-evaluation to other corpora and other published translations would show whether the pattern observed on TRACTOR appears under other programs and evaluation conditions.
